\begin{lemma}
\label{lemma-from-chow-to-K}
Let $X$ be a scheme locally of finite type over $(S, \delta)$
as in Situation \ref{situation-setup}. There is a canonical map
$$
\CH_k(X)
\longrightarrow
K_0(\textit{Coh}_{\leq k + 1}(X)/\textit{Coh}_{\leq k - 1}(X))
$$
induced by the map
$Z_k(X) \to K_0(\textit{Coh}_{\leq k}(X)/\textit{Coh}_{\leq k - 1}(X))$
from Lemma \ref{lemma-cycles-k-group}.
\end{lemma}

\begin{proof}
We have to show that an element $\alpha$ of $Z_k(X)$ which is rationally
equivalent to zero, is mapped to zero in
$K_0(\textit{Coh}_{\leq k + 1}(X)/\textit{Coh}_{\leq k - 1}(X))$.
Write $\alpha = \sum (i_j)_*\text{div}(f_j)$ as in
Definition \ref{definition-rational-equivalence}.
Observe that
$$
\pi = \coprod i_j : W = \coprod W_j \longrightarrow X
$$
is a finite morphism as each $i_j : W_j \to X$ is a closed immersion
and the family of $W_j$ is locally finite in $X$. Hence we may use
Lemma \ref{lemma-finite-cycles-k-group} to reduce to the case of $W$.
Since $W$ is a disjoint union of integral scheme, we reduce
to the case discussed in the next paragraph.

\medskip\noindent
Assume $X$ is integral of $\delta$-dimension $k + 1$.
Let $f$ be a nonzero rational function on $X$.
Let $\alpha = \text{div}(f)$. We have to show that
$\alpha$ is mapped to zero in
$K_0(\textit{Coh}_{\leq k + 1}(X)/\textit{Coh}_{\leq k - 1}(X))$.
Let $\mathcal{I} \subset \mathcal{O}_X$ be the ideal of denominators
of $f$, see Divisors, Definition
\ref{divisors-definition-regular-meromorphic-ideal-denominators}.
Then we have short exact sequences
$$
0 \to \mathcal{I} \to \mathcal{O}_X \to \mathcal{O}_X/\mathcal{I} \to 0
$$
and
$$
0 \to \mathcal{I} \xrightarrow{f} \mathcal{O}_X \to
\mathcal{O}_X/f\mathcal{I} \to 0
$$
See Divisors, Lemma
\ref{divisors-lemma-regular-meromorphic-ideal-denominators}.
We claim that
$$
[\mathcal{O}_X/\mathcal{I}]_k - [\mathcal{O}_X/f\mathcal{I}]_k =
\text{div}(f)
$$
The claim implies the element $\alpha = \text{div}(f)$ is represented by
$[\mathcal{O}_X/\mathcal{I}] - [\mathcal{O}_X/f\mathcal{I}]$
in $K_0(\textit{Coh}_{\leq k}(X)/\textit{Coh}_{\leq k - 1}(X))$.
Then the short exact sequences show that this element maps to
zero in $K_0(\textit{Coh}_{\leq k + 1}(X)/\textit{Coh}_{\leq k - 1}(X))$.

\medskip\noindent
To prove the claim, let $Z \subset X$ be an integral closed subscheme
of $\delta$-dimension $k$ and let $\xi \in Z$ be its generic point.
Then $I = \mathcal{I}_\xi \subset A = \mathcal{O}_{X, \xi}$
is an ideal such that $fI \subset A$. Now the coefficient of
$[Z]$ in $\text{div}(f)$ is $\text{ord}_A(f)$. (Of course as usual
we identify the function field of $X$ with the fraction field of $A$.)
On the other hand, the coefficient of $[Z]$ in
$[\mathcal{O}_X/\mathcal{I}] - [\mathcal{O}_X/f\mathcal{I}]$
is
$$
\text{length}_A(A/I) - \text{length}_A(A/fI)
$$
Using the distance function of
Algebra, Definition \ref{algebra-definition-distance}
we can rewrite this as
$$
d(A, I) - d(A, fI) = d(I, fI) = \text{ord}_A(f)
$$
The equalities hold by Algebra, Lemmas
\ref{algebra-lemma-properties-distance-function} and
\ref{algebra-lemma-order-vanishing-determinant}.
(Using these lemmas isn't necessary, but convenient.)
\end{proof}